\documentclass[11pt]{article}
\usepackage{amsmath,amssymb,amsthm,bm}
\usepackage[margin=1in]{geometry}
\usepackage{booktabs}
\usepackage{array}
\usepackage{enumitem}
\usepackage[hidelinks]{hyperref}

\theoremstyle{plain}
\newtheorem{proposition}{Proposition}
\newtheorem{corollary}{Corollary}
\theoremstyle{remark}
\newtheorem{remark}{Remark}

\newcommand{\E}{\mathbb{E}}
\newcommand{\pto}{\xrightarrow{p}}
\newcommand{\Xt}{X_\tau}
\newcolumntype{P}[1]{>{\raggedright\arraybackslash}p{#1}}

\title{\vspace{-3em}Scaling Ordinal Neighborhood--Wealth Estimation\\
to Cross-Country Zero-Shot Transfer\\[0.3em]
\large A Research Plan}
\author{}
\date{\today}

\begin{document}
\maketitle
\vspace{-2.5em}

\section{Objective and central thesis}
We have a building-level ordinal model that ranks structural wealth from
high-resolution aerial imagery, trained and validated on New York City
(2010--2024; out-of-sample Spearman $\rho=0.85$). The plan turns this into a
high-impact paper by changing the \emph{target domain} from a survey-rich setting
to settings where comparable small-area income data are stale, asset-only, or
absent. The thesis: \emph{a model trained only on U.S.\ imagery and labels
produces valid, temporally consistent neighborhood-wealth estimates in cities
with no contemporaneous local labels}, recovered by mapping the model's ordinal
scores onto a local distribution. This directly answers the ``why not just
download footprints and the ACS'' objection, which is unanswerable wherever such
data do not exist.

\section{Contributions, ranked}
\begin{enumerate}[leftmargin=1.4em,itemsep=0.15em]
\item \textbf{Zero-shot cross-country transfer (headline).} The U.S.-trained
ordinal model recovers neighborhood wealth in Mexico City and Vancouver with no
local training, validated against local ground truth.
\item \textbf{Identification of the estimand (methodological depth).} The score is
rank-consistent for \emph{location-expected wealth}, a kernel-smoother object; a
shrinkage corollary characterizes its bias and a testable implication makes the
building-level claim falsifiable.
\item \textbf{U.S.\ multi-city training on NAIP (infrastructure / fallback).} A
national, uniform training pipeline that both enables transfer and stands alone if
transfer underperforms.
\end{enumerate}

\section{Model and estimand (recap)}
The feature extractor is a ScaleMAE ViT-L backbone adapted with LoRA; a late-fusion
head emits a scalar score $R_b$ for the building tile $\Xt(b)$ (a $200\,$m window,
$\tau=100\,$m, subsampled to $224^2$). Training uses a RankNet cross-sectional loss
on within-year tract comparisons, a temporal stability penalty on structurally
unchanged twins, and a unit-variance regularizer; deployment maps ordinal scores to
a target city's income distribution by parametric (GB2) or empirical quantile
matching. The identified object is the \emph{location-expected wealth}
$m_\tau(b)=\E[\,w\mid \Xt(b)\,]$ (companion note for full statements):

\begin{proposition}[Locational identification]
The score is rank-consistent for location-expected wealth:
$R_b\pto\varphi(m_\tau(b))$ for a strictly increasing $\varphi$; i.e.\ it identifies
the ordering of the building's \emph{location}, not its idiosyncratic wealth.
\end{proposition}
\begin{corollary}[Shrinkage and bounds]
Writing $w_b=m_\tau(b)+e_b$, the score is a shrinkage estimator toward the local
mean with reliability $\rho_\tau\in[0,1]$ decreasing in $\tau$; building-level bias
$(1-\rho_\tau)(w_b-\bar w_i)$ vanishes for representative buildings and is largest
for outliers; $\tau$ is the bias--variance bandwidth.
\end{corollary}
\begin{remark}[Testable on NYC, no building-level labels needed]
Within-tract score dispersion should increase in observable within-tract
built-environment heterogeneity and vanish as it does---an external check on the
locational interpretation.
\end{remark}

\section{Data plan}
The effective model input is $200\,\mathrm{m}/224\,\mathrm{px}\approx 0.9\,$m/px, so
all three domains are brought to a common effective resolution; ScaleMAE absorbs
residual GSD differences. Imagery, centroids, construction timing, and validation
labels are sourced as follows.

\begin{table}[h]
\centering\footnotesize
\begin{tabular}{P{2.2cm} P{3.7cm} P{3.7cm} P{3.7cm}}
\toprule
 & \textbf{United States (train)} & \textbf{Mexico City (CDMX)} & \textbf{Vancouver} \\
\midrule
Imagery &
NAIP, 0.6--1\,m, $\sim$2009--present, free (National Map / USDA); leaf-on &
INEGI ortho, $\le 1$\,m to $<\!30$\,cm, free via geoportal; sub-metre basemap as fallback &
City open data, 7.5\,cm, 2015/2018/2020/2022, free \\[0.4em]
Footprints (centroids) &
Overture (GERS persistent IDs) &
Google Open Buildings v3 / Overture / Microsoft (single-epoch) &
City footprint layer (current $+$ dated vintages) \\[0.4em]
Construction timing (temporal) &
ACS B25034 ``Year Structure Built'' (tract); assessor data for showcase cities &
Google Open Buildings 2.5D \emph{Temporal}, annual 2016--2023, 4\,m raster $\to$ aggregate to AGEB &
Observed footprints 1999/2009/2015/current $+$ BC Assessment year built \\[0.4em]
Ground truth (labels / validation) &
ACS tract per-capita income (5-yr) &
INEGI AGEB census $+$ CONAPO marginalization index (2010/2020) &
StatCan dissemination-area income (2011/2016/2021) \\[0.4em]
Role &
Train $+$ internal OOS (NYC) &
External transfer: mid-income, harder &
External transfer: rich, clean \\
\bottomrule
\end{tabular}
\end{table}

\noindent\textbf{Key asymmetry.} Training needs only centroids and local income
labels---\emph{no} construction timing. Timing is required only for the event-study
and stability validation, which are tract-level; hence both transfer cities can be
validated cross-sectionally with single-epoch footprints, and temporally where
timing data exist (clean in Vancouver, AI-derived and tract-aggregated in CDMX).

\section{Experimental design (validation arc)}
\begin{enumerate}[leftmargin=1.4em,itemsep=0.2em]
\item \textbf{Training infrastructure.} Retrain the ordinal model on NAIP across a
national multi-city U.S.\ sample (ACS labels, Overture centroids). This diverse
distribution is what makes zero-shot transfer credible, and is the fallback result.
\item \textbf{Internal validity (have it).} NYC spatial-and-temporal holdout
(unseen tracts $\times$ 2016 holdout year, $\rho=0.85$).
\item \textbf{External validity (headline).} Zero-shot to Vancouver and CDMX;
quantile-map ordinal scores to local distributions; validate against StatCan DA
income and INEGI/CONAPO respectively.
\item \textbf{Temporal validation.} Vancouver: full Callaway--Sant'Anna event study
on new construction $+$ a building-level trajectory figure (observed footprints).
CDMX: tract-/AGEB-level event study using the AI temporal layer for
new-construction intensity (no building-level figure from the 4\,m raster).
\item \textbf{Identification result.} State the location-expected-wealth estimand,
the shrinkage corollary, and \emph{test} the dispersion implication on NYC.
\end{enumerate}

\section{Usefulness framing}
Do not claim these cities lack data---they do not. Claim the dimensions that are
scarce even there: (i) annual frequency vs.\ a decennial census; (ii) building-level
resolution finer than any administrative unit; (iii) a single comparable instrument
across heterogeneous data environments; and (iv) coverage at the truly data-void
margin. The CDMX result---run entirely on free, AI-generated temporal data with no
register---demonstrates the pipeline can operate where no administrative data exist,
which is the global generalization claim.

\section{Risks and mitigations}
\begin{itemize}[leftmargin=1.4em,itemsep=0.15em]
\item \textbf{Transfer failure on alien morphology} (informal/self-built housing in
CDMX). Mitigate: lead with Vancouver (morphologically close, clean ground truth) to
prove the mechanism; present CDMX as the harder, higher-value case; report partial
transfer ($\rho\!\approx\!0.6$--$0.7$ with zero local labels) as a positive.
\item \textbf{Domain shift US$\to$targets} (leaf-on/off, sensor, color). Mitigate:
train on NAIP so train/test effective resolution and sensor family align; rely on
photometric augmentation and the ordinal (distribution-free) design.
\item \textbf{AI temporal data quality in CDMX} (4\,m raster, inter-frame
registration drift, false positives). Mitigate: use it only for tract-level
treatment intensity, never building-level geometry.
\item \textbf{Cross-vintage footprint noise in Vancouver} (heterogeneous sources).
Mitigate: use BC Assessment year built for timing; treat snapshot differencing as a
cross-check.
\end{itemize}

\section{Target venues and deliverables}
Aim at the Jean (2016)/Yeh (2020)/Khachiyan (2022)/Zheng (2026) lineage: a
top general-science or economics-insights venue, where the data-scarce-transfer hook
plus a clean identification result is rewarded. Deliverables: (i) the retrained
multi-city NAIP model; (ii) zero-shot validation tables for Vancouver and CDMX;
(iii) temporal/event-study figures; (iv) the identification note and its NYC test.

\section{Indicative task sequence}
\begin{table}[h]
\centering\footnotesize
\begin{tabular}{P{0.8cm} P{12.5cm}}
\toprule
\textbf{T} & \textbf{Task} \\
\midrule
1 & Assemble NAIP $+$ Overture $+$ ACS national multi-city training set; retrain ordinal model. \\
2 & Reconfirm NYC internal holdout under the NAIP-trained model. \\
3 & Source Vancouver (7.5\,cm) and CDMX (INEGI/basemap) imagery; build centroid tiles. \\
4 & Zero-shot inference $+$ quantile mapping; validate vs.\ StatCan DA and INEGI/CONAPO. \\
5 & Vancouver event study $+$ building-level figure; CDMX AGEB-level event study via OB Temporal. \\
6 & Finalize identification proposition/corollary; run NYC dispersion test. \\
7 & Write-up; framing on frequency/resolution/comparability/data-void margin. \\
\bottomrule
\end{tabular}
\end{table}

\paragraph{Primary data sources.}
NAIP (USDA/USGS National Map); ACS B25034 and tract income (U.S.\ Census Bureau);
Overture Maps Buildings; Google Open Buildings v3 and 2.5D Temporal (CC-BY-4.0 /
ODbL); INEGI orthoimagery and AGEB census; CONAPO marginalization index; City of
Vancouver Open Data (orthophotos, footprints); BC Assessment; Statistics Canada
dissemination-area census.

\end{document}